\documentclass[11pt]{article}

\usepackage[letterpaper,margin=0.88in]{geometry}
\usepackage[T1]{fontenc}
\usepackage{lmodern}
\usepackage{microtype}
\usepackage{amsmath,amssymb,bm}
\usepackage{booktabs,tabularx,array,multirow}
\usepackage{graphicx}
\usepackage{float}
\usepackage{xcolor}
\usepackage{enumitem}
\usepackage{listings}
\usepackage{caption}
\usepackage{subcaption}
\usepackage{fancyhdr}
\usepackage{titlesec}
\usepackage[colorlinks=true,linkcolor=red!65!black,citecolor=red!65!black,urlcolor=red!65!black]{hyperref}

\definecolor{codeblue}{RGB}{35,80,140}
\definecolor{codegreen}{RGB}{25,110,65}
\definecolor{lightgray}{RGB}{246,246,246}
\definecolor{rulegray}{RGB}{185,185,185}

\setlength{\parindent}{0pt}
\setlength{\parskip}{0.48em}
\setlength{\headheight}{14pt}
\setlist[itemize]{leftmargin=1.5em,itemsep=0.25em,topsep=0.25em}
\setlist[enumerate]{leftmargin=1.7em,itemsep=0.25em,topsep=0.25em}
\captionsetup{font=small,labelfont=bf}
\titleformat{\section}{\Large\bfseries}{\thesection}{0.7em}{}
\titleformat{\subsection}{\large\bfseries}{\thesubsection}{0.7em}{}
\titleformat{\subsubsection}{\normalsize\bfseries}{\thesubsubsection}{0.7em}{}

\pagestyle{fancy}
\fancyhf{}
\lhead{VAE on MNIST: Reproduction Write-Up}
\rhead{Auto-Encoding Variational Bayes}
\cfoot{\thepage}
\renewcommand{\headrulewidth}{0.35pt}

\lstdefinestyle{pythonstyle}{
  language=Python,
  basicstyle=\ttfamily\footnotesize,
  keywordstyle=\color{codeblue}\bfseries,
  commentstyle=\color{codegreen},
  stringstyle=\color{red!65!black},
  backgroundcolor=\color{lightgray},
  frame=single,
  rulecolor=\color{rulegray},
  breaklines=true,
  columns=fullflexible,
  showstringspaces=false,
  aboveskip=0.5em,
  belowskip=0.5em
}

\newcommand{\R}{\mathbb{R}}
\newcommand{\E}{\mathbb{E}}
\newcommand{\KL}{D_{\mathrm{KL}}}
\newcommand{\code}[1]{\texttt{#1}}
\graphicspath{{figures/}}

\begin{document}

\begin{center}
  {\Large Reproduction Write-Up}\\[0.35em]
  {\LARGE\bfseries Auto-Encoding Variational Bayes on MNIST}\\[0.85em]
  {\large Based on Kingma and Welling, \emph{Auto-Encoding Variational Bayes}}\\[0.5em]
  {\normalsize Dataset: MNIST \quad Latent space: 20 dimensions \quad Hidden units: 500}\\[0.2em]
  {\normalsize Reproduction date: October 3, 2026}
\end{center}

\vspace{0.6em}
\section*{Start Here}

\begin{itemize}
  \item \textbf{Goal.} Reproduce the paper's MNIST variational auto-encoder (VAE) with a 20-dimensional latent space, including the evidence lower bound (ELBO), Gaussian reparameterization, stochastic optimization, reconstruction, prior sampling, t-SNE inspection, and latent interpolation.
  \item \textbf{Allowed software.} The model is implemented in PyTorch; \code{torchvision} is used only for the MNIST loader and image grids. The optimizer and automatic differentiation are provided by PyTorch.
  \item \textbf{Reproducibility.} The run fixes seed 42, stores all hyperparameters and metrics, and generates every figure in this write-up directly from the final checkpoint.
  \item \textbf{Scope.} This is a focused reproduction, not a claim of bit-for-bit equivalence. The paper's wake-sleep and Monte Carlo EM baselines are not repeated. All departures are listed in Section~\ref{sec:fidelity}.
\end{itemize}

\vspace{0.3em}
\section*{Deliverables}

\begin{itemize}
  \item \code{vae\_mnist\_reproduction\_writeup.pdf}: this report.
  \item \code{vae\_mnist\_training/}: runnable multi-file training project with documented class and function names.
  \item \code{vae\_mnist\_training.zip}: portable archive containing source, smoke test, trained checkpoint, configuration, metrics, history, and figures.
  \item \code{metrics.json} and \code{history.csv}: top-level copies of the final scalar results and per-epoch history.
\end{itemize}

\newpage
\section*{Reproduction Objectives}

After completing the experiment, the reader should be able to:

\begin{itemize}
  \item explain why direct maximum-likelihood learning is difficult when $p_\theta(\bm{x})=\int p_\theta(\bm{x},\bm{z})\,d\bm{z}$ is intractable;
  \item derive the ELBO and identify its reconstruction and KL-divergence terms;
  \item explain why $\bm{z}=\bm{\mu}+\bm{\sigma}\odot\bm{\epsilon}$ permits low-variance pathwise gradients;
  \item map each mathematical object to its PyTorch implementation and tensor shape;
  \item reproduce the reported metrics and visual diagnostics from a clean environment.
\end{itemize}

\section*{Checklist}

\begin{tabularx}{\textwidth}{@{}p{0.055\textwidth}X@{}}
  $\square$ & Install Python, PyTorch, torchvision, NumPy, and Matplotlib. \\
  $\square$ & Install packages and run \code{train.py} with the command in Appendix~\ref{app:run}. \\
  $\square$ & Confirm that MNIST contains 60,000 training and 10,000 test images. \\
  $\square$ & Confirm that both posterior heads output shape $N\times 20$. \\
  $\square$ & Confirm that reconstruction and KL losses are finite and the KL term is non-negative. \\
  $\square$ & Confirm that train and test negative ELBO decrease without a widening generalization gap. \\
  $\square$ & Inspect reconstructions, prior samples, the t-SNE embedding, and latent interpolations. \\
  $\square$ & Compare the final metrics with Table~\ref{tab:results}. \\
\end{tabularx}

\vfill
\begin{center}
\fbox{\parbox{0.9\textwidth}{\textbf{Reading note.} The paper maximizes the ELBO. The implementation minimizes its negative. Therefore, a lower reported ``negative ELBO'' is better.}}
\end{center}

\newpage
\tableofcontents
\newpage

\section{Introduction}

Kingma and Welling propose the stochastic gradient variational Bayes (SGVB) estimator and the auto-encoding variational Bayes (AEVB) algorithm for directed latent-variable models with intractable posteriors~\cite{kingma2013}. The key idea is to learn two networks jointly:

\begin{itemize}
  \item a \textbf{probabilistic encoder} $q_\phi(\bm{z}\mid\bm{x})$ that approximates the posterior over latent variables; and
  \item a \textbf{probabilistic decoder} $p_\theta(\bm{x}\mid\bm{z})$ that models how observations are generated from a latent code.
\end{itemize}

Unlike a deterministic auto-encoder, the VAE is a normalized generative model. The latent prior constrains the code space, while a likelihood model supplies a statistically meaningful reconstruction term. Reparameterization moves the sampling noise into an auxiliary variable independent of $\phi$, allowing gradients to pass through the sample.

MNIST is a suitable reproduction target because the paper uses it directly and its pixels can be modeled with Bernoulli likelihoods. A preliminary two-dimensional run was easy to visualize but produced blurry reconstructions and limited sample diversity. The primary experiment therefore uses $J=20$, one of the latent sizes studied in the paper. Its geometry is inspected with t-SNE and controlled interpolations rather than forcing the generative model itself to remain two-dimensional. The experiment asks three concrete questions:

\begin{enumerate}
  \item Does stochastic optimization improve the held-out ELBO?
  \item Does the learned posterior reconstruct previously unseen digits?
  \item Does sampling and moving continuously through the prior produce coherent digit-like images?
\end{enumerate}

\section{Problem Setup and Notation}

Let $\bm{x}\in[0,1]^{784}$ be a flattened $28\times 28$ MNIST image and $\bm{z}\in\R^J$ a continuous latent vector. The model factorizes as
\begin{equation}
  p_\theta(\bm{x},\bm{z})=p(\bm{z})p_\theta(\bm{x}\mid\bm{z}),
  \qquad p(\bm{z})=\mathcal{N}(\bm{0},\bm{I}).
  \label{eq:joint}
\end{equation}
The true posterior $p_\theta(\bm{z}\mid\bm{x})$ is intractable because evaluating it requires the marginal likelihood
\begin{equation}
  p_\theta(\bm{x})=\int p(\bm{z})p_\theta(\bm{x}\mid\bm{z})\,d\bm{z}.
\end{equation}
We therefore introduce a diagonal Gaussian approximation:
\begin{equation}
  q_\phi(\bm{z}\mid\bm{x})=
  \mathcal{N}\!\left(\bm{z};\bm{\mu}_\phi(\bm{x}),
  \operatorname{diag}\!\left(\bm{\sigma}_\phi^2(\bm{x})\right)\right).
  \label{eq:q}
\end{equation}

\begin{table}[H]
\centering
\caption{Notation used throughout the reproduction.}
\label{tab:notation}
\begin{tabularx}{0.94\textwidth}{@{}l l l X@{}}
\toprule
Symbol & Code name & Shape & Meaning \\
\midrule
$N$ & batch size & scalar & Number of images in a minibatch; $N=100$. \\
$\bm{x}$ & \code{x} & $N\times784$ & Flattened input images. \\
$\bm{h}_e$ & encoder hidden & $N\times500$ & Tanh encoder representation. \\
$\bm{\mu},\log\bm{\sigma}^2$ & \code{mu}, \code{logvar} & $N\times J$ & Approximate posterior parameters. \\
$\bm{\epsilon}$ & noise & $N\times J$ & Standard-normal auxiliary noise. \\
$\bm{z}$ & latent code & $N\times J$ & Reparameterized posterior sample. \\
$\bm{a}$ & \code{logits} & $N\times784$ & Bernoulli logits produced by the decoder. \\
$\phi$ & encoder parameters & -- & Weights of encoder, mean head, and log-variance head. \\
$\theta$ & decoder parameters & -- & Weights of decoder hidden and output layers. \\
\bottomrule
\end{tabularx}
\end{table}

\section{Variational Objective}

\subsection{Evidence lower bound}

Insert $q_\phi(\bm{z}\mid\bm{x})$ into the marginal log-likelihood and apply Jensen's inequality:
\begin{align}
\log p_\theta(\bm{x})
&=\log\int q_\phi(\bm{z}\mid\bm{x})
  \frac{p_\theta(\bm{x},\bm{z})}{q_\phi(\bm{z}\mid\bm{x})}\,d\bm{z} \\
&\geq \E_{q_\phi(\bm{z}\mid\bm{x})}
\left[\log p_\theta(\bm{x},\bm{z})-\log q_\phi(\bm{z}\mid\bm{x})\right] \\
&=\E_{q_\phi}\left[\log p_\theta(\bm{x}\mid\bm{z})\right]
-\KL\!\left(q_\phi(\bm{z}\mid\bm{x})\,\|\,p(\bm{z})\right)
\equiv \mathcal{L}(\theta,\phi;\bm{x}).
\label{eq:elbo}
\end{align}
The first term rewards accurate decoding. The KL term regularizes each approximate posterior toward the standard-normal prior, making prior sampling meaningful.

\subsection{Reparameterization trick}

Sampling $\bm{z}\sim q_\phi(\bm{z}\mid\bm{x})$ naively places a stochastic node between the loss and $\phi$. For a Gaussian posterior, write
\begin{equation}
  \bm{\epsilon}\sim\mathcal{N}(\bm{0},\bm{I}),
  \qquad
  \bm{z}=\bm{\mu}_\phi(\bm{x})+
  \exp\!\left(\tfrac12\log\bm{\sigma}_\phi^2(\bm{x})\right)\odot\bm{\epsilon}.
  \label{eq:reparam}
\end{equation}
The random variable $\bm{\epsilon}$ is independent of $\phi$. Once sampled, the right-hand side is an ordinary differentiable computation, so automatic differentiation yields a pathwise gradient. As in the paper, the reproduction uses one latent sample per datapoint ($L=1$).

\subsection{Closed-form KL and implemented loss}

For the diagonal Gaussian posterior and standard-normal prior,
\begin{equation}
  \KL(q_\phi\|p)=
  -\frac12\sum_{j=1}^{J}
  \left(1+\log\sigma_j^2-\mu_j^2-\sigma_j^2\right).
  \label{eq:kl}
\end{equation}
MNIST pixels are modeled with independent Bernoulli variables. If $a_d$ is a decoder logit, the per-image negative ELBO minimized in code is
\begin{equation}
  \mathcal{J}(\bm{x})=
  \underbrace{-\sum_{d=1}^{784}\left[x_d\log \sigma(a_d)+(1-x_d)\log(1-\sigma(a_d))\right]}_{\text{binary cross-entropy / negative log-likelihood}}
  +\underbrace{\KL(q_\phi\|p)}_{\text{latent regularizer}}.
  \label{eq:loss}
\end{equation}
For numerical stability, the implementation passes logits directly to
\path{binary_cross_entropy_with_logits}.

\section{Reproduction Design}

\subsection{Network architecture}

The architecture follows the paper's one-hidden-layer MLP example. Both hidden layers use tanh, and all linear weights are initialized from $\mathcal{N}(0,0.01^2)$.

\begin{table}[H]
\centering
\caption{Model architecture for the final $J=20$ experiment.}
\label{tab:architecture}
\begin{tabularx}{0.92\textwidth}{@{}l l l X@{}}
\toprule
Stage & Operation & Output shape & Role \\
\midrule
Input & Flatten $28\times28$ image & $N\times784$ & Observation vector. \\
Encoder & $\tanh(\bm{x}W_e^\top+\bm{b}_e)$ & $N\times500$ & Shared inference features. \\
Posterior heads & Linear$(500,20)$ twice & $N\times20$ each & Produce $\bm{\mu}$ and $\log\bm{\sigma}^2$. \\
Sampling & Equation~\eqref{eq:reparam} & $N\times20$ & Differentiable stochastic code. \\
Decoder & $\tanh(\bm{z}W_d^\top+\bm{b}_d)$ & $N\times500$ & Generative hidden representation. \\
Output & Linear$(500,784)$ & $N\times784$ & Bernoulli logits. \\
\bottomrule
\end{tabularx}
\end{table}

\subsection{Dataset and optimization}

\begin{table}[H]
\centering
\caption{Final experimental configuration.}
\label{tab:hyperparameters}
\begin{tabularx}{0.84\textwidth}{@{}l X@{}}
\toprule
Setting & Value \\
\midrule
Dataset & MNIST: 60,000 training and 10,000 test images \\
Input treatment & Scale to $[0,1]$; dynamic Bernoulli binarization during training \\
Posterior / prior & Diagonal Gaussian / standard normal \\
Latent dimension & 20 \\
Hidden units & 500 in encoder and 500 in decoder \\
Minibatch / samples & 100 images / one $\bm{z}$ sample per image \\
Optimizer & Adagrad, learning rate 0.01 \\
Epochs & 50 (30,000 parameter updates) \\
Initialization & Independent $\mathcal{N}(0,0.01^2)$ weights, zero biases \\
Random seed & 42 \\
Evaluation & One posterior sample; unbinarized $[0,1]$ test targets \\
Device & Apple Metal Performance Shaders (MPS) \\
\bottomrule
\end{tabularx}
\end{table}

\subsection{Fidelity to the original paper}
\label{sec:fidelity}

\begin{table}[H]
\centering
\caption{What is reproduced and what is intentionally simplified.}
\begin{tabularx}{\textwidth}{@{}p{0.23\textwidth}X X@{}}
\toprule
Component & Paper & This reproduction \\
\midrule
Core estimator & AEVB/SGVB with analytic Gaussian KL and $L=1$ & Same \\
MNIST model & Bernoulli decoder, Gaussian encoder, one 500-unit hidden layer & Same \\
Batch / optimizer & Batch 100; Adagrad selected from $\{0.01,0.02,0.1\}$ & Batch 100; fixed Adagrad 0.01 \\
Latent dimensions & Multiple dimensions for bound comparisons, including 20; 2-D used for visualization & One 20-D primary model for reconstruction and generation quality; t-SNE/interpolation for visualization \\
Training budget & Reported by evaluated examples, extending to many millions & 50 epochs = 3.0 million evaluated training examples \\
Baselines & Wake-sleep and, in one setting, Monte Carlo EM & Omitted \\
Weight decay & Small MAP-style weight decay, coefficient unspecified & Omitted rather than inventing a coefficient \\
Evaluation & Variational bound and separate marginal-likelihood study & Negative ELBO plus visual diagnostics \\
\bottomrule
\end{tabularx}
\end{table}

\section{Implementation}

\subsection{Training files and function interfaces}

The deliverable is a runnable project rather than a notebook. The map below gives the actual file name, the callable signature to look for, and the returned value. Bodies not essential to the interface are abbreviated with comments; the complete implementations are in \code{vae\_mnist\_training/}.

\begin{lstlisting}[style=pythonstyle,caption={Configuration, data, model, and loss interfaces.}]
# vae_mnist/config.py
@dataclass
class TrainConfig:
    data_dir: str = "work/data"
    output_dir: str = "work/experiment_z20"
    epochs: int = 50
    batch_size: int = 100
    hidden_dim: int = 500
    latent_dim: int = 20
    learning_rate: float = 0.01
    optimizer: str = "adagrad"
    seed: int = 42
    dynamic_binarization: bool = True

# vae_mnist/data.py
def build_mnist_loaders(config: TrainConfig):
    # Download MNIST and create seeded loaders.
    return train_loader, test_loader
\end{lstlisting}

\begin{lstlisting}[style=pythonstyle,caption={Model and objective interfaces.}]
# vae_mnist/model.py
class VariationalAutoencoder(nn.Module):
    def reset_parameters(self) -> None:
        # Initializes linear weights from N(0, 0.01^2).

    def encode(self, x):
        return mu, logvar

    @staticmethod
    def reparameterize(mu, logvar):
        std = torch.exp(0.5 * logvar)
        return mu + std * torch.randn_like(std)

    def decode_logits(self, z):
        return logits

    def forward(self, x):
        # Compose encode, reparameterize, and decode_logits.
        return logits, mu, logvar

# vae_mnist/losses.py
def negative_elbo(logits, target, mu, logvar):
    # All three tensors are sums over the current minibatch.
    return nelbo, reconstruction, kl
\end{lstlisting}

\begin{lstlisting}[style=pythonstyle,caption={Optimization, utility, and entry-point interfaces.}]
# vae_mnist/engine.py
def run_epoch(model, loader, device, optimizer=None,
              dynamic_binarization=False):
    # optimizer=None selects evaluation mode.
    return {"nelbo": nelbo_per_image,
            "reconstruction": recon_per_image,
            "kl": kl_per_image}

def fit(model, train_loader, test_loader, optimizer,
        device, config):
    # One dictionary per epoch, with train/test ELBO terms.
    return history

# vae_mnist/utils.py
def seed_everything(seed: int) -> None:
    # Seeds Python, NumPy, PyTorch, and MPS.

def choose_device():
    # Selects CUDA first, then MPS, then CPU.
    return device

def save_json(value, path: Path) -> None:
    # Writes a reproducible metadata file.
\end{lstlisting}

\begin{lstlisting}[style=pythonstyle,caption={Checkpoint and training entry-point interfaces.}]
# vae_mnist/checkpoint.py
def save_checkpoint(path, model, optimizer, config, metrics):
    # Writes model/optimizer states, config, and metrics.
    return checkpoint_path

def load_checkpoint(path, device):
    # Restores weights and switches the model to eval mode.
    return model, config, checkpoint

# train.py
def build_argument_parser():
    return parser

def config_from_args(args):
    return TrainConfig(...)

def build_optimizer(model, config):
    return optimizer  # Adagrad by default; Adam is optional.

def main() -> None:
    # Calls fit(), then save_checkpoint(..., "checkpoint.pt").
\end{lstlisting}

\begin{lstlisting}[style=pythonstyle,caption={Checkpoint inference used to generate Figures 2, 4, and 5.}]
# vae_mnist/inference.py
def encode(model, images, device):
    return mu, logvar

def reconstruct(model, images, device):
    # Decode posterior means; shape is [N, 1, 28, 28].
    return reconstructions

def interpolate(model, start_images,
                end_images, device, steps=10):
    # Shape is [number_of_pairs, steps, 1, 28, 28].
    return decoded_paths

def generate(model, num_samples, device):
    # Draw z ~ N(0,I); shape is [num_samples, 1, 28, 28].
    return generated_images

# vae_mnist/visualization.py
def save_reconstructions(model, test_loader, device, figure_dir):
    return None  # Writes reconstructions.png (Figure 2).

def save_latent_interpolations(model, test_loader,
                               device, figure_dir):
    return None  # Writes latent_interpolations.png (Figure 4).

def save_prior_samples(model, device, figure_dir):
    return None  # Writes prior_samples.png (Figure 5).

# infer.py
def main() -> None:
    model, config, checkpoint = load_checkpoint(path, device)
    # --task reconstruct, interpolate, generate, or all
    return None

# tests/test_smoke.py
class VaeSmokeTest(unittest.TestCase):
    def test_checkpoint_and_inference_interfaces(self) -> None:
        # Checks checkpoint round-trip and every inference shape.
\end{lstlisting}

The \code{pretrained/} directory contains \code{checkpoint.pt}, the saved configuration, metrics, history, and generated figures. A training run starts at \code{train.py::main}. General inference starts at \code{infer.py::main}; \code{--task all} regenerates Figures~\ref{fig:reconstruction}, \ref{fig:interpolation}, and \ref{fig:samples} from the checkpoint without optimization. The broader \code{render\_checkpoint.py::main} entry point additionally recreates the training curves and t-SNE figure.

\subsection{Training algorithm}

The stochastic training loop is the minibatch AEVB algorithm specialized to MNIST:

\begin{enumerate}
  \item Sample a shuffled minibatch of 100 images and dynamically binarize its pixels.
  \item Encode each image to $(\bm{\mu},\log\bm{\sigma}^2)$.
  \item Draw one $\bm{\epsilon}\sim\mathcal{N}(\bm{0},\bm{I})$ per image and construct $\bm{z}$.
  \item Decode $\bm{z}$ into 784 Bernoulli logits.
  \item Sum binary cross-entropy and analytic KL over the batch, divide by batch size, and backpropagate.
  \item Update $\theta$ and $\phi$ jointly with Adagrad.
\end{enumerate}

Dynamic binarization resamples each training pixel as Bernoulli$(x_d)$ on every visit. This matches the discrete Bernoulli observation model while using the original grayscale intensity as the sampling probability. Test targets remain in $[0,1]$ to make the held-out curve less noisy.

\subsection{Reproducibility controls}

The \path{train.py} entry point seeds Python, NumPy, and PyTorch; uses a seeded data-loader generator; and calls \path{save_checkpoint} to store the model state, optimizer state, configuration, and metrics. The matching \path{load_checkpoint} function reconstructs the model in evaluation mode. The run is self-contained: MNIST is downloaded automatically if absent. The separate \path{infer.py} and \path{render_checkpoint.py} entry points recreate outputs without repeating training.

\section{Results}

\subsection{Quantitative results}

\begin{table}[H]
\centering
\caption{Held-out comparison after 50 epochs. Lower negative ELBO and bits/dimension are better.}
\label{tab:results}
\begin{tabular}{@{}lrrl@{}}
\toprule
Metric & Preliminary $J=2$ & Final $J=20$ & Change \\
\midrule
Test negative ELBO & 166.065 & \textbf{112.599} & $-53.466$ nats/image \\
Test reconstruction term & 161.175 & \textbf{87.900} & Sharper reconstructions \\
Test KL term & 4.891 & \textbf{24.698} & More encoded information \\
Bits/dimension & 0.3056 & \textbf{0.2072} & 32.2\% reduction \\
Training examples evaluated & 3,000,000 & 3,000,000 & Equal data budget \\
\bottomrule
\end{tabular}
\end{table}

The 20-D model improves the primary metric by 53.466 nats/image relative to the preliminary 2-D model under the same epoch and minibatch budget. Figure~\ref{fig:curves} shows that its test negative ELBO falls from 171.4 after the first epoch to 112.6 at epoch 50. The final train--test gap is only 2.34 nats/image. The KL contribution grows to 24.7 nats/image instead of collapsing, confirming that the additional latent coordinates carry useful information.

\begin{figure}[H]
\centering
\includegraphics[width=0.98\textwidth]{training_curves.png}
\caption{Optimization history and decomposition of the held-out negative ELBO.}
\label{fig:curves}
\end{figure}

\newpage
\subsection{Reconstruction}

For reconstruction, the test image is encoded and the posterior mean $\bm{\mu}(\bm{x})$ is decoded. Using the mean removes sampling noise and tests whether the learned representation retains digit identity. Figure~\ref{fig:reconstruction} shows sharp posterior-mean reconstructions across all twelve examples. Compared with $J=2$, thin strokes, loops, and digit-specific geometry are substantially better preserved. It is generated directly by \code{infer.py --task reconstruct}, which calls \code{load\_checkpoint}, \code{reconstruct}, and \code{save\_reconstructions}.

\begin{figure}[H]
\centering
\includegraphics[width=0.98\textwidth]{reconstructions.png}
\caption{Top row: unseen MNIST inputs. Bottom row: posterior-mean reconstructions.}
\label{fig:reconstruction}
\end{figure}

\subsection{Latent representation}

Because the learned representation is 20-dimensional, plotting two raw coordinates would be arbitrary. Figure~\ref{fig:tsne} embeds 5,000 posterior means with t-SNE using perplexity 30, PCA initialization, 1,000 iterations, and random seed 42. The well-separated digit clusters indicate that the unsupervised representation captures class-related structure. Local neighborhoods are meaningful, but the global distances and cluster sizes should not be interpreted as metric properties of the original 20-D space. Labels are used only after training for coloring.

\begin{figure}[H]
\centering
\includegraphics[width=0.72\textwidth]{latent_tsne.png}
\caption{Fixed-seed t-SNE embedding of 5,000 twenty-dimensional posterior means (perplexity 30). Labels are used only for coloring after training.}
\label{fig:tsne}
\end{figure}

\newpage
\subsection{Generative behavior}

Figure~\ref{fig:interpolation} linearly interpolates between posterior means for five pairs of held-out digits. Every row changes gradually from its source to its destination: $0\rightarrow1$, $2\rightarrow3$, $4\rightarrow5$, $6\rightarrow7$, and $8\rightarrow9$. Smooth intermediate images show that the 20-D latent representation has learned coherent local geometry even though it cannot be displayed directly as a plane. The figure is generated by \code{infer.py --task interpolate}, which calls \code{interpolate} and \code{save\_latent\_interpolations}.

\begin{figure}[H]
\centering
\includegraphics[width=0.98\textwidth]{latent_interpolations.png}
\caption{Ten-step posterior-mean interpolations between five pairs of held-out digit classes.}
\label{fig:interpolation}
\end{figure}

Independent samples $\bm{z}\sim\mathcal{N}(\bm{0},\bm{I})$ in Figure~\ref{fig:samples} produce substantially more varied digit-like images than the 2-D model. The \code{generate} task in \code{infer.py} calls \code{generate} and \path{save_prior_samples}. Ambiguous samples remain because the decoder factorizes pixels as independent Bernoulli variables and the run is shorter than the paper's longest training horizons.

\begin{figure}[H]
\centering
\includegraphics[width=0.68\textwidth]{prior_samples.png}
\caption{One hundred independent samples from the learned decoder with $\bm{z}\sim\mathcal{N}(\bm{0},\bm{I})$.}
\label{fig:samples}
\end{figure}

\section{Validation and Interpretation}

\subsection{Sanity checks}

\begin{table}[H]
\centering
\caption{Checks used to decide whether the reproduction was successful.}
\begin{tabularx}{\textwidth}{@{}p{0.28\textwidth}X p{0.13\textwidth}@{}}
\toprule
Check & Evidence & Status \\
\midrule
Objective improves & Test negative ELBO decreases by 58.85 nats/image from epoch 1 to 50. & Pass \\
Generalization is stable & Final train--test gap is 2.34 nats/image. & Pass \\
Posterior is active & Final KL is 24.698 nats/image, not approximately zero. & Pass \\
Reconstructions retain identity & Posterior-mean outputs preserve strokes, loops, and digit identity. & Pass \\
Latent structure is meaningful & Fixed-seed t-SNE forms digit-related clusters; interpolations are smooth. & Pass \\
Prior generates recognizable data & Prior samples are varied and predominantly digit-like. & Pass \\
Exact paper comparison & No identical training horizon or baseline was available. & Not claimed \\
\bottomrule
\end{tabularx}
\end{table}

\subsection{What the experiment demonstrates}

The reproduction supports the paper's central mechanism. A single differentiable objective trains the inference and generative networks jointly. The analytic KL term supplies a principled regularizer, while reparameterization permits ordinary backpropagation through latent samples. The learned prior is not merely a compression device: sampling from it produces new observations, and continuous paths in latent space decode to continuous changes in image space.

The experiment also makes the reconstruction--regularization trade-off visible. The preliminary two-dimensional code was convenient to plot directly but compressed away stroke details. Increasing to $J=20$ sharply improves reconstruction and sample diversity. Direct visualization is recovered after training with t-SNE for neighborhood structure and explicit latent interpolations for continuity; neither visualization changes the learned model.

\section{Conclusion and Limitations}

This write-up reproduces the paper's MNIST VAE with the same probabilistic ingredients and closely matched architectural choices. After 50 epochs, the final 20-D model reaches a held-out negative ELBO of 112.599 nats/image, 53.466 nats/image better than the preliminary two-dimensional model. It reconstructs unseen digits sharply, forms class-related neighborhoods in a fixed-seed t-SNE embedding, supports smooth semantic interpolations, and generates diverse digit-like samples from the standard-normal prior.

The main limitations are deliberate. First, the final study reports one 20-D run and uses the 2-D run only as a diagnostic comparison; it is not a multi-seed latent-size ablation. Second, the wake-sleep and Monte Carlo EM baselines are omitted. Third, training covers 3 million examples, whereas the paper plots substantially longer horizons. Fourth, the paper mentions a small weight-decay term without giving its coefficient, so this reproduction omits it rather than choosing an undocumented value. Fifth, t-SNE preserves local neighborhoods but can distort global distances. Finally, the one-sample ELBO is a lower bound and should not be read as an exact marginal log-likelihood.

A stronger follow-up would train $J\in\{2,5,10,20\}$ under an equal compute budget, report multiple seeds, add importance-weighted likelihood estimates, and measure sample quality with a classifier-based precision/recall analysis. Those extensions are useful scientific comparisons but are not required to validate the AEVB mechanism reproduced here.

\appendix
\section{How to Run the Reproduction}
\label{app:run}

From the \code{vae\_mnist\_training/} project directory:

\begin{lstlisting}[style=pythonstyle]
python3 -m pip install -r requirements.txt
python3 train.py \
  --epochs 50 \
  --batch-size 100 \
  --hidden-dim 500 \
  --latent-dim 20 \
  --learning-rate 0.01 \
  --optimizer adagrad \
  --seed 42 \
  --data-dir ../../work/data \
  --output-dir ../../work/experiment_z20_project
\end{lstlisting}

After training, load the saved checkpoint and regenerate Figure~\ref{fig:reconstruction}, Figure~\ref{fig:interpolation}, and Figure~\ref{fig:samples} without optimization:

\begin{lstlisting}[style=pythonstyle]
python3 infer.py \
  --checkpoint pretrained/checkpoint.pt \
  --data-dir ../../work/data \
  --output-dir ../../work/inference_figures \
  --task all \
  --seed 42
\end{lstlisting}

Replace \code{all} with \code{reconstruct}, \code{interpolate}, or \code{generate} to produce only Figure~\ref{fig:reconstruction}, Figure~\ref{fig:interpolation}, or Figure~\ref{fig:samples}, respectively.

Required packages:

\begin{lstlisting}[style=pythonstyle]
python3 -m pip install torch==2.9.1 torchvision==0.24.1 \
  numpy matplotlib pillow scikit-learn
\end{lstlisting}

The output directory contains \code{checkpoint.pt}, \code{config.json}, \code{metrics.json}, \code{history.csv}, and five PNG figures. On CUDA or MPS hardware the script selects an accelerator automatically; otherwise it runs on CPU. Validate shapes and gradient flow before a full run with:

\begin{lstlisting}[style=pythonstyle]
python3 -m unittest tests/test_smoke.py
\end{lstlisting}

To regenerate figures from the delivered checkpoint, use \path{render_checkpoint.py} as documented in the project README.

\section{Loss-to-Code Map}

\begin{table}[H]
\centering
\caption{Direct correspondence between the derivation and implementation.}
\begin{tabularx}{\textwidth}{@{}p{0.28\textwidth}p{0.31\textwidth}X@{}}
\toprule
Mathematical quantity & PyTorch expression & Note \\
\midrule
$\bm{\sigma}$ & \code{torch.exp(0.5 * logvar)} & Log-variance is numerically convenient. \\
$\bm{z}$ & \code{mu + std * randn\_like(std)} & Equation~\eqref{eq:reparam}. \\
$-\log p_\theta(\bm{x}\mid\bm{z})$ & \code{binary\_cross\_entropy}\newline\code{\_with\_logits} & Sums over pixels and batch. \\
$\KL(q_\phi\|p)$ & \code{-0.5 * sum(1 + logvar -}\newline\code{mu**2 - exp(logvar))} & Exact diagonal-Gaussian result. \\
$-\mathcal{L}$ & \code{reconstruction + kl} & Minimized by the optimizer. \\
\bottomrule
\end{tabularx}
\end{table}

\section{References}

\begin{thebibliography}{9}

\bibitem{kingma2013}
D. P. Kingma and M. Welling.
\newblock Auto-Encoding Variational Bayes.
\newblock \emph{arXiv:1312.6114}, 2013.

\bibitem{lecun1998}
Y. LeCun, L. Bottou, Y. Bengio, and P. Haffner.
\newblock Gradient-based learning applied to document recognition.
\newblock \emph{Proceedings of the IEEE}, 86(11):2278--2324, 1998.

\bibitem{duchi2011}
J. Duchi, E. Hazan, and Y. Singer.
\newblock Adaptive subgradient methods for online learning and stochastic optimization.
\newblock \emph{Journal of Machine Learning Research}, 12:2121--2159, 2011.

\bibitem{pytorch}
A. Paszke et al.
\newblock PyTorch: An imperative style, high-performance deep learning library.
\newblock \emph{Advances in Neural Information Processing Systems}, 2019.

\end{thebibliography}

\end{document}
